% Figura: punto fuente, punto campo y vector que los une.
\begin{figure}[ht]
\centering
\tdplotsetmaincoords{70}{118}
\begin{tikzpicture}[tdplot_main_coords, scale=1.55]
  \def\sx{1.5}\def\sy{-0.6}\def\sz{2.0}   % punto fuente r'
  \def\fx{1.4}\def\fy{2.4}\def\fz{0.4}   % punto campo r
  \draw[eje] (0,0,0) -- (3.0,0,0) node[below left] {\small $x$};
  \draw[eje] (0,0,0) -- (0,3.1,0) node[right] {\small $y$};
  \draw[eje] (0,0,0) -- (0,0,3.1) node[above] {\small $z$};
  % vectores de posición
  \draw[->, line width=1.4pt, azulacero] (0,0,0) -- (\sx,\sy,\sz);
  \draw[->, line width=1.4pt, azulprof] (0,0,0) -- (\fx,\fy,\fz);
  % vector relativo R = r - r'
  \draw[->, line width=1.8pt, rojoteo] (\sx,\sy,\sz) -- (\fx,\fy,\fz);
  % versor u_R en el punto campo
  \pgfmathsetmacro{\Rx}{\fx-\sx}\pgfmathsetmacro{\Ry}{\fy-\sy}\pgfmathsetmacro{\Rz}{\fz-\sz}
  \pgfmathsetmacro{\Rn}{sqrt(\Rx*\Rx+\Ry*\Ry+\Rz*\Rz)}
  \draw[->, line width=1.6pt, verdeteo] (\fx,\fy,\fz) -- ++({0.95*\Rx/\Rn},{0.95*\Ry/\Rn},{0.95*\Rz/\Rn});
  \fill[azulacero] (\sx,\sy,\sz) circle (1.8pt);
  \fill[azulprof] (\fx,\fy,\fz) circle (1.8pt);
  \fill (0,0,0) circle (1.2pt);
  % etiquetas
  \node[azulacero, font=\normalsize] at (0.35,-0.55,1.0) {$\vc{r}\,'$};
  \node[azulprof, font=\normalsize] at (0.55,1.35,-0.15) {$\vc{r}$};
  \node[rojoteo, font=\normalsize, fill=white, fill opacity=0.8, text opacity=1, inner sep=1pt]
     at (1.55,1.15,1.5) {$\vc{R}=\vc{r}-\vc{r}\,'$};
  \node[verdeteo, font=\normalsize] at ({\fx+1.3*\Rx/\Rn},{\fy+1.3*\Ry/\Rn},{\fz+1.3*\Rz/\Rn}) {$\uR$};
  \node[font=\small, text=black!80, anchor=east] at (\sx+0.05,\sy-0.05,\sz+0.1) {fuente};
  \node[font=\small, text=black!80, anchor=north] at (\fx,\fy,\fz-0.12) {campo};
\end{tikzpicture}
\caption{Funciones de la posición relativa. El vector $\vc{r}\,'$ (azul claro)
localiza el \emph{punto fuente}, $\vc{r}$ (azul oscuro) el \emph{punto campo}, y
$\vc{R}=\vc{r}-\vc{r}\,'$ (rojo) apunta de la fuente al punto donde se evalúa
el campo; $\uR=\vc{R}/R$ (verde) es su versor.}
\label{fig:posicion-relativa}
\end{figure}
